\chapter{Weak Derivatives}
\label{ch:iii20}

Weak derivatives extend differentiation to $L^1_{\mathrm{loc}}$ functions by testing against compactly supported smooth functions. They are the bridge from distribution theory to Sobolev spaces and weak PDE formulations.

\section*{Learning goals}
The reader should be able to state the central definitions precisely, prove the structural results, audit the hypotheses in examples, and use the chapter's estimates in later analysis.

\section*{Conceptual roadmap}
\[
\boxed{\text{structure}\;\longrightarrow\;\text{estimate}\;\longrightarrow\;\text{limit or representation}\;\longrightarrow\;\text{application}.}
\]

\section{Weak derivative definition}
A function $g$ is the weak derivative $D_jf$ if $\int fD_j\phi=-\int g\phi$ for every test function.

\section{Agreement with classical derivatives}
Classically differentiable functions have the same weak derivative.

\section{Piecewise smooth functions}
Jumps usually prevent the weak derivative from being represented by an ordinary function because the distributional derivative contains Dirac masses.

\section{Uniqueness}
Weak derivatives are unique up to almost-everywhere equality.

\section{Higher weak derivatives}
Iterating the definition yields multi-index weak derivatives.

\section{Locality}
Weak derivatives depend only on local almost-everywhere behavior.

\section{Integration by parts}
The weak derivative definition is precisely the boundary-free integration-by-parts identity.

\section{Approximation}
Mollification converts weak derivatives into classical derivatives and commutes with differentiation on interior regions.

\section{Core structural results}

\begin{theorem}[Uniqueness of weak derivative]\label{thm:iii20-01}
If $g,h\in L^1_{\mathrm{loc}}$ both satisfy the weak derivative identity for $f$, then $g=h$ almost everywhere.
\begin{proof}
Their difference integrates to zero against every test function; the fundamental lemma of distributions implies the associated regular distribution is zero, hence the functions agree a.e.
\end{proof}
\end{theorem}

\begin{theorem}[Classical implies weak]\label{thm:iii20-02}
If $f\in C^1$, then its weak derivative equals its classical derivative.
\begin{proof}
Integrate by parts against compactly supported test functions; boundary terms vanish.
\end{proof}
\end{theorem}

\begin{theorem}[Mollifier commutation]\label{thm:iii20-03}
For locally integrable $f$ with weak derivative $D_jf$, $D_j(f*\rho_\varepsilon)=(D_jf)*\rho_\varepsilon$ where defined.
\begin{proof}
Differentiate the smooth convolution and transfer the derivative from the mollifier to $f$ using the weak derivative identity.
\end{proof}
\end{theorem}

\section{Worked examples}

\begin{example}[Weak derivative definition diagnostic]\label{ex:iii20-worked-01}
Give a rigorous worked analysis of Weak derivative definition in the setting of this chapter. State the decisive definition or hypothesis and derive the central conclusion. A function $g$ is the weak derivative $D_jf$ if $\int fD_j\phi=-\int g\phi$ for every test function. The decisive point is to use the chapter definition at the correct countable, norm, transform, or distributional level rather than rely on pointwise intuition alone.
\end{example}

\begin{example}[Agreement with classical derivatives diagnostic]\label{ex:iii20-worked-02}
Give a rigorous worked analysis of Agreement with classical derivatives in the setting of this chapter. State the decisive definition or hypothesis and derive the central conclusion. Classically differentiable functions have the same weak derivative. The decisive point is to use the chapter definition at the correct countable, norm, transform, or distributional level rather than rely on pointwise intuition alone.
\end{example}

\begin{example}[Piecewise smooth functions diagnostic]\label{ex:iii20-worked-03}
Give a rigorous worked analysis of Piecewise smooth functions in the setting of this chapter. State the decisive definition or hypothesis and derive the central conclusion. Jumps usually prevent the weak derivative from being represented by an ordinary function because the distributional derivative contains Dirac masses. The decisive point is to use the chapter definition at the correct countable, norm, transform, or distributional level rather than rely on pointwise intuition alone.
\end{example}

\begin{example}[Uniqueness diagnostic]\label{ex:iii20-worked-04}
Give a rigorous worked analysis of Uniqueness in the setting of this chapter. State the decisive definition or hypothesis and derive the central conclusion. Weak derivatives are unique up to almost-everywhere equality. The decisive point is to use the chapter definition at the correct countable, norm, transform, or distributional level rather than rely on pointwise intuition alone.
\end{example}

\section{Solved dossiers}
Each dossier is a canonical solved problem. The provenance ledger records which retained corpus rules guided its topic and which dossiers were added freshly to complete the chapter.

\begin{problem}[Weak derivative definition diagnostic]\label{prob:iii20-dossier-01}
Give a rigorous worked analysis of Weak derivative definition in the setting of this chapter. State the decisive definition or hypothesis and derive the central conclusion.
\end{problem}
\begin{solution}
A function $g$ is the weak derivative $D_jf$ if $\int fD_j\phi=-\int g\phi$ for every test function. The decisive point is to use the chapter definition at the correct countable, norm, transform, or distributional level rather than rely on pointwise intuition alone.
\end{solution}

\begin{problem}[Agreement with classical derivatives diagnostic]\label{prob:iii20-dossier-02}
Give a rigorous worked analysis of Agreement with classical derivatives in the setting of this chapter. State the decisive definition or hypothesis and derive the central conclusion.
\end{problem}
\begin{solution}
Classically differentiable functions have the same weak derivative. The decisive point is to use the chapter definition at the correct countable, norm, transform, or distributional level rather than rely on pointwise intuition alone.
\end{solution}

\begin{problem}[Piecewise smooth functions diagnostic]\label{prob:iii20-dossier-03}
Give a rigorous worked analysis of Piecewise smooth functions in the setting of this chapter. State the decisive definition or hypothesis and derive the central conclusion.
\end{problem}
\begin{solution}
Jumps usually prevent the weak derivative from being represented by an ordinary function because the distributional derivative contains Dirac masses. The decisive point is to use the chapter definition at the correct countable, norm, transform, or distributional level rather than rely on pointwise intuition alone.
\end{solution}

\begin{problem}[Uniqueness diagnostic]\label{prob:iii20-dossier-04}
Give a rigorous worked analysis of Uniqueness in the setting of this chapter. State the decisive definition or hypothesis and derive the central conclusion.
\end{problem}
\begin{solution}
Weak derivatives are unique up to almost-everywhere equality. The decisive point is to use the chapter definition at the correct countable, norm, transform, or distributional level rather than rely on pointwise intuition alone.
\end{solution}

\begin{problem}[Higher weak derivatives diagnostic]\label{prob:iii20-dossier-05}
Give a rigorous worked analysis of Higher weak derivatives in the setting of this chapter. State the decisive definition or hypothesis and derive the central conclusion.
\end{problem}
\begin{solution}
Iterating the definition yields multi-index weak derivatives. The decisive point is to use the chapter definition at the correct countable, norm, transform, or distributional level rather than rely on pointwise intuition alone.
\end{solution}

\begin{problem}[Locality diagnostic]\label{prob:iii20-dossier-06}
Give a rigorous worked analysis of Locality in the setting of this chapter. State the decisive definition or hypothesis and derive the central conclusion.
\end{problem}
\begin{solution}
Weak derivatives depend only on local almost-everywhere behavior. The decisive point is to use the chapter definition at the correct countable, norm, transform, or distributional level rather than rely on pointwise intuition alone.
\end{solution}

\begin{problem}[Integration by parts diagnostic]\label{prob:iii20-dossier-07}
Give a rigorous worked analysis of Integration by parts in the setting of this chapter. State the decisive definition or hypothesis and derive the central conclusion.
\end{problem}
\begin{solution}
The weak derivative definition is precisely the boundary-free integration-by-parts identity. The decisive point is to use the chapter definition at the correct countable, norm, transform, or distributional level rather than rely on pointwise intuition alone.
\end{solution}

\begin{problem}[Approximation diagnostic]\label{prob:iii20-dossier-08}
Give a rigorous worked analysis of Approximation in the setting of this chapter. State the decisive definition or hypothesis and derive the central conclusion.
\end{problem}
\begin{solution}
Mollification converts weak derivatives into classical derivatives and commutes with differentiation on interior regions. The decisive point is to use the chapter definition at the correct countable, norm, transform, or distributional level rather than rely on pointwise intuition alone.
\end{solution}

\begin{problem}[Uniqueness of weak derivative proof dossier]\label{prob:iii20-dossier-09}
Prove the following structural statement and identify the step where its main hypothesis is used: If $g,h\in L^1_{\mathrm{loc}}$ both satisfy the weak derivative identity for $f$, then $g=h$ almost everywhere.
\end{problem}
\begin{solution}
Their difference integrates to zero against every test function; the fundamental lemma of distributions implies the associated regular distribution is zero, hence the functions agree a.e. This proof also explains why the stated hypothesis is structural rather than cosmetic.
\end{solution}

\begin{problem}[Classical implies weak proof dossier]\label{prob:iii20-dossier-10}
Prove the following structural statement and identify the step where its main hypothesis is used: If $f\in C^1$, then its weak derivative equals its classical derivative.
\end{problem}
\begin{solution}
Integrate by parts against compactly supported test functions; boundary terms vanish. This proof also explains why the stated hypothesis is structural rather than cosmetic.
\end{solution}

\begin{problem}[Mollifier commutation proof dossier]\label{prob:iii20-dossier-11}
Prove the following structural statement and identify the step where its main hypothesis is used: For locally integrable $f$ with weak derivative $D_jf$, $D_j(f*\rho_\varepsilon)=(D_jf)*\rho_\varepsilon$ where defined.
\end{problem}
\begin{solution}
Differentiate the smooth convolution and transfer the derivative from the mollifier to $f$ using the weak derivative identity. This proof also explains why the stated hypothesis is structural rather than cosmetic.
\end{solution}

\begin{problem}[Chapter synthesis]\label{prob:iii20-dossier-12}
Connect the main definitions and structural theorems of this chapter into one proof strategy. Explain what object is controlled, what estimate or closure principle is used, and what conclusion becomes available.
\end{problem}
\begin{solution}
Weak derivatives extend differentiation to $L^1_{\mathrm{loc}}$ functions by testing against compactly supported smooth functions. They are the bridge from distribution theory to Sobolev spaces and weak PDE formulations. The recurring strategy is to encode the object in the chapter's natural measurable, normed, Fourier, or distributional structure, establish the controlling estimate, and then pass to the desired limit or representation.
\end{solution}

\section{Exercises with complete solutions}
The shorter exercise layer remains separate from the solved dossiers.

\begin{exercise}\label{exr:iii20-01}
State and apply the central principle behind Weak derivative definition. Give one immediate consequence in the setting of this chapter.
\end{exercise}
\begin{hint}
Start from the precise definition or estimate in the corresponding section.
\end{hint}
\begin{solution}
A function $g$ is the weak derivative $D_jf$ if $\int fD_j\phi=-\int g\phi$ for every test function. As an immediate consequence, the same principle can be inserted into later convergence, approximation, transform, or weak-form arguments whenever its hypotheses are verified.
\end{solution}

\begin{exercise}\label{exr:iii20-02}
State and apply the central principle behind Agreement with classical derivatives. Give one immediate consequence in the setting of this chapter.
\end{exercise}
\begin{hint}
Start from the precise definition or estimate in the corresponding section.
\end{hint}
\begin{solution}
Classically differentiable functions have the same weak derivative. As an immediate consequence, the same principle can be inserted into later convergence, approximation, transform, or weak-form arguments whenever its hypotheses are verified.
\end{solution}

\begin{exercise}\label{exr:iii20-03}
State and apply the central principle behind Piecewise smooth functions. Give one immediate consequence in the setting of this chapter.
\end{exercise}
\begin{hint}
Start from the precise definition or estimate in the corresponding section.
\end{hint}
\begin{solution}
Jumps usually prevent the weak derivative from being represented by an ordinary function because the distributional derivative contains Dirac masses. As an immediate consequence, the same principle can be inserted into later convergence, approximation, transform, or weak-form arguments whenever its hypotheses are verified.
\end{solution}

\begin{exercise}\label{exr:iii20-04}
State and apply the central principle behind Uniqueness. Give one immediate consequence in the setting of this chapter.
\end{exercise}
\begin{hint}
Start from the precise definition or estimate in the corresponding section.
\end{hint}
\begin{solution}
Weak derivatives are unique up to almost-everywhere equality. As an immediate consequence, the same principle can be inserted into later convergence, approximation, transform, or weak-form arguments whenever its hypotheses are verified.
\end{solution}

\begin{exercise}\label{exr:iii20-05}
State and apply the central principle behind Higher weak derivatives. Give one immediate consequence in the setting of this chapter.
\end{exercise}
\begin{hint}
Start from the precise definition or estimate in the corresponding section.
\end{hint}
\begin{solution}
Iterating the definition yields multi-index weak derivatives. As an immediate consequence, the same principle can be inserted into later convergence, approximation, transform, or weak-form arguments whenever its hypotheses are verified.
\end{solution}

\begin{exercise}\label{exr:iii20-06}
State and apply the central principle behind Locality. Give one immediate consequence in the setting of this chapter.
\end{exercise}
\begin{hint}
Start from the precise definition or estimate in the corresponding section.
\end{hint}
\begin{solution}
Weak derivatives depend only on local almost-everywhere behavior. As an immediate consequence, the same principle can be inserted into later convergence, approximation, transform, or weak-form arguments whenever its hypotheses are verified.
\end{solution}

\begin{exercise}\label{exr:iii20-07}
State and apply the central principle behind Integration by parts. Give one immediate consequence in the setting of this chapter.
\end{exercise}
\begin{hint}
Start from the precise definition or estimate in the corresponding section.
\end{hint}
\begin{solution}
The weak derivative definition is precisely the boundary-free integration-by-parts identity. As an immediate consequence, the same principle can be inserted into later convergence, approximation, transform, or weak-form arguments whenever its hypotheses are verified.
\end{solution}

\begin{exercise}\label{exr:iii20-08}
State and apply the central principle behind Approximation. Give one immediate consequence in the setting of this chapter.
\end{exercise}
\begin{hint}
Start from the precise definition or estimate in the corresponding section.
\end{hint}
\begin{solution}
Mollification converts weak derivatives into classical derivatives and commutes with differentiation on interior regions. As an immediate consequence, the same principle can be inserted into later convergence, approximation, transform, or weak-form arguments whenever its hypotheses are verified.
\end{solution}

\section*{Chapter summary}
The chapter now supplies a canonical theorem-and-dossier layer for subsequent measure, Fourier, distributional, Sobolev, and PDE arguments.
